%% Delete this \nocite command invocation to make the references section only list out the bibitems that are actually cited.
\nocite{*}




\subsection{Surfaces}

\begin{definition} \label{definition:surface_as_a_manifold}
A \hldef{surface} is a two-\CrefAndHyperrefIfExist{definition:dimension_of_a_topological_space}{dimensional} \CrefAndHyperrefIfExist{definition:topological_manifold}{topological manifold}.
\end{definition}

\begin{proposition} \label{proposition:surface_is_orientable_if_and_only_if_it_does_not_contain_a_subset_homeomorphic_to_a_mobius_strip}
    \TODO{mobius strip}
A surface is orientable if and only if it does not contain any subset homeomorphic to a Möbius strip.
\end{proposition}

\subsection{Euler characteristic and genus}

\begin{definition} \label{definition:euler_characteristic_of_a_topological_space}
    \TODO{betti numbers}
Let $X$ be a topological space and let $H_n(X; \mathbb{Q})$ be the $n$-th singular homology group of $X$ with coefficients in the rational numbers $\mathbb{Q}$. If the total dimension of the singular homology is finite, the \hldef{Euler characteristic of $X$}, denoted \hl{$\chi(X)$}, is defined by the alternating sum of the Betti numbers:
$$\chi(X) = \sum_{n=0}^{\infty} (-1)^n \dim_{\mathbb{Q}} H_n(X; \mathbb{Q})$$
\end{definition}

\begin{definition} \label{definition:genus_of_a_topological_space}
Let $X$ be a topological space. 
One might denote by \hl{$g = g(X)$} the unique number in $\frac{1}{2} \bbZ$ defined by
$$g = \frac{2-\chi(X)}{2}$$
or equivalently 
$$\chi(X) = 2 - 2g$$
where $\chi(X)$ is the Euler characteristic of $X$ defined via singular homology.

In the case that $X$ is a \CrefAndHyperrefIfExist{definition:surface_as_a_manifold}{surface}, the quantity $g = g(X)$ above may be called the \hldef{genus of the surface}, especially when the surface is compact, connected, and orientable. For compact, connected, non-orientable surfaces, the genus may be called the \hldef{"Non-orientable genus"} or \hldef{"demigenus"} for emphasis. In the case that $X$ is a disconnected surface, the above notion of genus should not be confused with the notion of total genus. 

If the surface is compact, connected, and orientable, then its genus is a non-negative integer. If the surface is compact, connected, and non-orientable, then its genus is a non-negative half integer.
\end{definition}


\begin{theorem} \label{theorem:compact_connected_orientable_surfaces_are_homeomorphic_if_and_only_if_they_have_the_same_genus}
Let $S$ and $S'$ be two \CrefAndHyperrefIfExist{definition:compact_topological_space}{compact}, \CrefAndHyperrefIfExist{definition:disconnected_and_connected_topological_spaces}{connected}, orientable surfaces. $S$ and $S'$ have the same \CrefAndHyperref{definition:homeomorphism_of_topological_spaces}{homeomorphism class} if and only if they have the same \CrefAndHyperref{definition:genus_of_a_topological_space}{genus} $g$. Consequently, for every non-negative integer $g$, there exists a unique compact, connected, orientable surface \hl{$S_g$}, called by names such as the \hldef{closed orientable surface of genus $g$}, the \hldef{surface of genus $g$}, or the \hldef{$g$-fold torus},  up to \CrefAndHyperrefIfExist{definition:homeomorphism_of_topological_spaces}{homeomorphism}.
\end{theorem}

\section{The Burau representations}

\subsection{Braid groups}

\subsubsection{Configuration spaces}

\begin{definition} \label{definition:ordered_configuration_space_of_n_points_in_a_topological_space}
Let $X$ be a \CrefAndHyperrefIfExist{definition:category_of_topological_spaces}{topological space}. 
The \hldef{ordered configuration space of $n$ points in $X$}, denoted \hl{$\operatorname{Conf}_n(X)$}, is the subspace of the $n$-fold \CrefAndHyperrefIfExist{definition:cartesian_product_of_topological_spaces}{product space} $X^n$ defined by
$$\operatorname{Conf}_n(X) = \{ (x_1, \dots, x_n) \in X^n : x_i \neq x_j \text{ for } i \neq j \}.$$
Equivalently, $\operatorname{Conf}_n(X) = X^n \setminus \Delta$, where $\Delta = \{ (x_1, \dots, x_n) : \exists i \neq j \text{ s.t. } x_i = x_j \}$ is the fat diagonal.
\TODO{fat diagonal}
\end{definition}
\begin{definition} \label{definition:unordered_configuration_space_of_n_points_in_a_topological_space}
Let $X$ be a \CrefAndHyperrefIfExist{definition:category_of_topological_spaces}{topological space}.
The \hldef{unordered configuration space of $n$ points in $X$}, denoted \hl{$\operatorname{UConf}_n(X)$}, is the \CrefAndHyperrefIfExist{definition:quotient_of_a_topological_space_with_respect_to_an_equivalence_relations}{quotient space}
$$\operatorname{UConf}_n(X) = \operatorname{Conf}_n(X) / S_n$$
where the \CrefAndHyperrefIfExist{definition:symmetric_group_on_a_set}{symmetric group} $S_n$ \CrefAndHyperrefIfExist{definition:left_and_right_group_actions_on_a_set}{acts on} $\operatorname{Conf}_n(X)$ by permutation of coordinates:
$$\sigma \cdot (x_1, \dots, x_n) = (x_{\sigma(1)}, \dots, x_{\sigma(n)})$$
for $\sigma \in S_n$. The points in $\operatorname{UConf}_n(X)$ are viewed as $n$-element subsets of $X$.
\end{definition}
\begin{definition} \label{definition:artin_braid_group}
    \TODO{presentation of a group}
The \hldef{(Artin) braid group} \hl{$B_n$} is the group with the presentation given by generators $\sigma_1, \dots, \sigma_{n-1}$ subject to the relations:
\begin{itemize}
    \item $\sigma_i \sigma_j = \sigma_j \sigma_i$ for $|i - j| \geq 2$
    \item $\sigma_i \sigma_{i+1} \sigma_i = \sigma_{i+1} \sigma_i \sigma_{i+1}$ for $1 \leq i \leq n-2$
\end{itemize}
An element of $B_n$ may be called a \hldef{braid}. The generators $\sigma_1,\ldots, \sigma_{n-1}$ may be called the \hldef{Artin generators}, \hldef{elementary braids}, \hldef{simple braids}, or \hldef{elementary crossings}.

Equivalently, $B_n$ can be characterized as the \CrefAndHyperrefIfExist{definition:homotopy_groups_of_a_pointed_topological_space}{fundamental group} $\pi_1(\operatorname{UConf}_n(\mathbb{D}, c_0)$ of the \CrefAndHyperrefIfExist{definition:unordered_configuration_space_of_n_points_in_a_topological_space}{unordered configuration space} $\operatorname{UConf}_n(\mathbb{D})$ of the (unit) disk $\mathbb{D} \subseteq \bbR^2$ based at any point $c_0 \in \operatorname{UConf}_n(\mathbb{D})$.
\end{definition}
\begin{proposition} \label{proposition:artin_braid_group_is_isomorphic_to_mapping_class_group_of_homeomorphisms_of_disk_with_n_punctured_points_fixing_the_boundary_and_permuting_the_punctures}
    Let $\bbD \subseteq \bbR^2$ be a (closed unit) disk and let $x_1,\ldots,x_n \in \bbD$ be $n$ distinct points in the interior of $\bbD$.
There is a canonical isomorphism
$$B_n \cong \text{Mod}(\bbD, \partial \mathbb{D}, \{x_1,\ldots,x_n\})$$
(\Cref{definition:artin_braid_group}) where $\text{Mod}(\bbD, \partial \bbD, \{x_1,\ldots,x_n\})$ is the \CrefAndHyperref{definition:mapping_class_group_of_a_connected_oriented_topological_manifold_with_or_without_boundary_relative_to_subsets}{mapping class group} of homeomorphisms of $\bbD$ fixing the boundary $\partial \mathbb{D}$ pointwise and permuting the punctures.
\end{proposition}
\begin{definition} \label{definition:pure_braid_group}
The \hldef{pure braid group}, denoted \hl{$P_n$}, is the \CrefAndHyperrefIfExist{definition:kernel_image_of_group_homomorphism}{kernel} of the natural \CrefAndHyperrefIfExist{definition:injective_surjective_bijective_map_of_sets}{surjective} \CrefAndHyperrefIfExist{definition:group_homomorphism}{homomorphism} $\pi: B_n \to S_n$ (\Cref{definition:artin_braid_group}) onto the \CrefAndHyperrefIfExist{definition:symmetric_group_on_a_set}{symmetric group} $S_n$, which maps a braid to the \CrefAndHyperrefIfExist{definition:symmetric_group_on_a_set}{permutation} of its endpoints.

Note in particular that $P_n$ is a subset of $B_n$, and hence all elements of $P_n$ are \CrefAndHyperrefIfExist{definition:artin_braid_group}{braids}.

Topologically, $P_n$ is canonically isomorphic to the \CrefAndHyperrefIfExist{definition:homotopy_groups_of_a_pointed_topological_space}{fundamental group} $\pi_1(\operatorname{Conf}_n(\bbD), c_0)$ of the \CrefAndHyperrefIfExist{definition:ordered_configuration_space_of_n_points_in_a_topological_space}{ordered configuration space} $\operatorname{Conf}_n(\bbD)$ of $n$ points in the disk $\bbD \subseteq \bbR^2$.
\end{definition}

Configuration spaces can also be constructed for schemes.

\begin{definition} \label{definition:ordered_configuration_space_of_n_points_in_a_scheme_over_a_scheme}
    Let $S$ be a \CrefAndHyperrefIfExist{definition:scheme}{scheme} and let $X$ be a \CrefAndHyperrefIfExist{definition:scheme_over_a_scheme}{scheme over $S$}. 
    The \hldef{ordered configuration space of $n$ points in $X$ (over $S$)} refers to the $n$-fold \CrefAndHyperrefIfExist{definition:cartesian_product_of_two_objects_in_a_category_over_an_object}{fiber product} $X^n = X \times_S X \times_S \cdots \times_S X$.

    Note that it \CrefAndHyperrefIfExist{definition:representable_functor_on_a_locally_small_category}{represents} the \CrefAndHyperrefIfExist{definition:functor_between_categories}{functor} $\Sch/S \to \Sets: T \mapsto \Hom_{\Sch/S}(T, X)^n$\CrefIfExists{definition:scheme_over_a_scheme}. 
\end{definition}



\begin{definition} \label{definition:ordered_configuration_space_of_n_distinct_points_in_a_scheme_over_a_scheme}
    Let $S$ be a \CrefAndHyperrefIfExist{definition:scheme}{scheme} and let $X$ be a \CrefAndHyperrefIfExist{definition:scheme_over_a_scheme}{scheme over $S$}. 
    The \hldef{configuration space of $n$ distinct points in $X$ (over $S$)} refers to the scheme 
    $$\operatorname{Conf}_n(X) = X^n \setminus \bigcup_{1 \leq i < j \leq n} \Delta_{ij}$$
    \CrefIfExists{definition:ordered_configuration_space_of_n_points_in_a_scheme_over_a_scheme}
    where $\Delta_{ij} \subseteq X^n = X \times_S X \times_S \cdots \times_S X$ refers to the pairwise diagonal at the $i$ and $j$th coordinates --- it may be constructed by first taking the diagonal $\Delta \subseteq X_i \times_S X_j = X \times_S X$ where $X_i$ and $X_j$ are $i$th and $j$th copies of $X$ and by letting $\Delta_{ij}$ be the product of $\Delta$ with the copies of $X$ except the $i$th and $j$th ones.
    \TODO{comment on how it represents tuples of points}

    \TextIfExists{definition:unordered_configuration_space_of_n_distinct_points_in_a_scheme_over_a_scheme}{
        This configuration space should not be confused with the \CrefAndHyperrefIfExist{definition:unordered_configuration_space_of_n_distinct_points_in_a_scheme_over_a_scheme}{unordered configuration space of $n$ distinct points in $X$ over $S$}, which some might also denote by $\operatorname{Conf}_n(X)$
    }
\end{definition}

\begin{definition} \label{definition:unordered_configuration_space_of_n_distinct_points_in_a_scheme_over_a_scheme}
    Let $S$ be a \CrefAndHyperrefIfExist{definition:scheme}{scheme} and let $X$ be a \CrefAndHyperrefIfExist{definition:scheme_over_a_scheme}{scheme over $S$}. 
    The \hldef{unordered configuration space of $n$ distinct points in $X$ (over $S$)} refers to the scheme 
    $$\hlin{\operatorname{UConf}_n(X) = \operatorname{Conf}_n(X) / S_n}$$
    \CrefIfExists{definition:ordered_configuration_space_of_n_distinct_points_in_a_scheme_over_a_scheme}\CrefIfExists{definition:symmetric_group_on_a_set} where $\operatorname{Conf}_n(X)$ is the ordered configuration space of $n$ distinct points in $X$ and $S_n$ acts on $\operatorname{Conf}_n(X)$ by permuting coordinates.

    This unordered configuration space may be denoted by \hl{$\operatorname{Conf}_n(X)$}, which may overlap with notation for the ordered configuration space.

    \TODO{comment on how it represents tuples of points}
\end{definition}

\begin{definition} \label{definition:unordered_configuration_space_of_colored_distinct_points_in_a_scheme_over_a_scheme}
    Let $S$ be a \CrefAndHyperrefIfExist{definition:scheme}{scheme} and let $X$ be a \CrefAndHyperrefIfExist{definition:scheme_over_a_scheme}{scheme over $S$}. 
    For integers $n_1,\ldots,n_v \geq 1$ and $n = \sum_{i=1}^v n_i$, the \hldef{unordered configuration space of $n_1,\ldots,n_v$ distinct colored points in $X$ (over $S$)} refers to the scheme 
    $$\hlin{\operatorname{UConf}_{n_1,\ldots,n_v}(X) = \operatorname{Conf}_n(X) / \prod_{i=1}^v S_{n_i}}$$
    \CrefIfExists{definition:ordered_configuration_space_of_n_distinct_points_in_a_scheme_over_a_scheme}\CrefIfExists{definition:product_of_groups}\CrefIfExists{definition:symmetric_group_on_a_set} where $S_{n_i}$ acts by permuting the $n_i$ consecutive coordinates in the range $[n_1+\cdots + n_{i-1} + 1, n_1 + \cdots + n_i]$. 
    \TODO{comment on how it represents tuples of points}

    This configuration space may also be denoted as $\operatorname{UConf}_{n_1,\ldots,n_v}(X)$.
\end{definition}





\subsection{Defining the Burau representation}

\subsubsection{Orientation by singular homology}

\begin{definition} \label{definition:local_orientation_of_a_topological_manifold_with_or_without_boundary_at_an_interior_point_using_local_singular_homology}
    \TODO{generator}
Let $M$ be a \CrefAndHyperref{definition:topological_manifold}{topological $n$-manifold (with or without boundary)}. 

A \hldef{local orientation of $M$ at a point $x \in \operatorname{int} M$} is a choice of generator $\mu_x$ of the local homology group $H_n(M, M \setminus \{x\}; \mathbb{Z}) \cong \mathbb{Z}$ \CrefIfExists{definition:singular_homology_and_relative_singular_homology_of_topological_space_with_coefficients_in_an_abelian_group}.
\end{definition}
\begin{definition} \label{definition:orientation_of_a_topological_manifold_with_or_without_boundary_using_singular_homology}
Let $M$ be a \CrefAndHyperref{definition:topological_manifold}{topological $n$-manifold (with or without boundary)}.

An \hldef{orientation of $M$} is a function $x \mapsto \mu_x$ assigning a \CrefAndHyperref{definition:local_orientation_of_a_topological_manifold_with_or_without_boundary_at_an_interior_point_using_local_singular_homology}{local orientation} to each point $x \in \operatorname{int} M$ \CrefIfExists{definition:interior_and_boundary_of_topological_space} such that for every $x \in \operatorname{int} M$, there exists an open neighborhood $U \subset \operatorname{int} M$ and a class $\alpha \in H_n(M, M \setminus U; \mathbb{Z})$ such that for all $y \in U$, the map induced by the inclusion $j_y: (M, M \setminus U) \to (M, M \setminus \{y\})$ sends $\alpha$ to $\mu_y$.

If $M$ is a topological manifold with or without boundary for which an orientation exists, then $M$ is said to be \hldef{orientable}.

A topological manifold $M$ with or without boundary equipped with an orientation is said to be an \hldef{oriented manifold}. 
\end{definition}
\begin{definition} \label{definition:orientation_preserving_homeomorphism_between_oriented_topological_manifolds_with_or_without_boundary}
Let $M$ and $N$ be \CrefAndHyperref{definition:topological_manifold}{topological $n$-manifolds (with or without boundary)}, equipped with \CrefAndHyperref{definition:orientation_of_a_topological_manifold_with_or_without_boundary_using_singular_homology}{orientations} $\mu^M$ and $\mu^N$ respectively.

A \CrefAndHyperrefIfExist{definition:homeomorphism_of_topological_spaces}{homeomorphism} $f: M \to N$ is \hldef{orientation-preserving} if for every $x \in \operatorname{int} M$, the induced isomorphism on local \CrefAndHyperref{definition:singular_homology_and_relative_singular_homology_of_topological_space_with_coefficients_in_an_abelian_group}{homology}
$$f_*: H_n(M, M \setminus \{x\}; \mathbb{Z}) \to H_n(N, N \setminus \{f(x)\}; \mathbb{Z})$$
maps the chosen generator $\mu^M_x$ to the chosen generator $\mu^N_{f(x)}$.
\end{definition}

\begin{definition} \label{definition:orientation_preserving_homeomorphism_group_of_a_topological_manifold_with_or_without_boundary}
Let $M$ be a \CrefAndHyperrefIfExist{definition:topological_manifold}{topological $n$-manifold (with or without boundary)}.
\TODO{homeomorphism group}
The \hldef{orientation-preserving homeomorphism group of $M$}, denoted \hl{$\text{Homeo}^{+}(M)$}, is the subgroup of the homeomorphism group $\text{Homeo}(M)$ consisting of all \CrefAndHyperref{definition:orientation_preserving_homeomorphism_between_oriented_topological_manifolds_with_or_without_boundary}{orientation-preserving homeomorphisms} $M \to M$.
\end{definition}


\begin{definition} \label{definition:orientation_preserving_homeomorphism_group_fixing_boundary_of_an_oriented_topological_manifold_with_or_without_boundary}
Let $M$ be an \CrefAndHyperref{definition:orientation_of_a_topological_manifold_with_or_without_boundary_using_singular_homology}{oriented} \CrefAndHyperref{definition:topological_manifold}{topological $n$-manifold with or without boundary}.
The \hldef{group of orientation-preserving homeomorphisms fixing the boundary}, denoted \hl{$\text{Homeo}^{+}(M, \partial M)$}, is the subgroup of $\text{Homeo}^{+}(M)$ (\Cref{definition:orientation_preserving_homeomorphism_group_of_a_topological_manifold_with_or_without_boundary}) defined by
$$\text{Homeo}^{+}(M, \partial M) = \{f \in \text{Homeo}^{+}(M) : f(x) = x \text{ for all } x \in \partial M\}.$$
Note that this group is simply $\operatorname{Homeo}^+(M, \partial M, \emptyset)$ in the notation of \Cref{notation:homeomorphism_groups_of_a_topological_space_manifolds}.
\end{definition}


\begin{theorem} \label{theorem:homeomorphism_from_a_compact_connected_oriented_topological_manifold_to_itself_is_orientation_preserving_iff_relative_singular_homology_map_is_identity}
Let $M$ be a \CrefAndHyperrefIfExist{definition:compact_topological_space}{compact}, \CrefAndHyperrefIfExist{definition:disconnected_and_connected_topological_spaces}{connected}, \CrefAndHyperrefIfExist{definition:orientation_of_a_topological_manifold_with_or_without_boundary_using_singular_homology}{oriented} \CrefAndHyperrefIfExist{topologicadefinition:topological_manifold}{topological $n$-manifold}, which may be with or without boundary. A \CrefAndHyperrefIfExist{definition:homeomorphism_of_topological_spaces}{homeomorphism} $f: M \to M$ is \CrefAndHyperref{definition:orientation_preserving_homeomorphism_between_oriented_topological_manifolds_with_or_without_boundary}{orientation-preserving} if and only if the induced map $f_*: H_n(M, \partial M; \mathbb{Z}) \to H_n(M, \partial M; \mathbb{Z})$ on the \CrefAndHyperrefIfExist{definition:singular_homology_and_relative_singular_homology_of_topological_space_with_coefficients_in_an_abelian_group}{relative singular homology}  is the identity.
\end{theorem}



\subsection{Isotopy between topological spaces and between homeomorphisms of topological spaces}

\begin{definition} \label{definition:isotopy_between_topological_spaces}
Let $M$ and $N$ be \CrefAndHyperrefIfExist{definition:topological_spaces}{topological spaces}. 

An \hldef{isotopy from $M$ to $N$} is a \CrefAndHyperref{definition:homotopy_of_maps_of_topological_spaces_relative_to_a_subset}{homotopy} $H: M \times [0, 1] \to N$ such that for each $t \in [0, 1]$, the map $h_t: M \to N$ defined by $h_t(x) = H(x, t)$ is a \CrefAndHyperrefIfExist{definition:homeomorphism_of_topological_spaces}{homeomorphism}.

In practice, this terminology is used essentially exclusively in the case that $M$ and $N$ are \CrefAndHyperref{definition:topological_manifold}{topological manifolds with or without boundary}.

\end{definition}
\begin{definition} \label{definition:isotopic_homeomorphisms_between_topological_spaces}
Let $M$ and $N$ be \CrefAndHyperrefIfExist{definition:topological_spaces}{topological spaces}. 

Two \CrefAndHyperrefIfExist{definition:homeomorphism_of_topological_spaces}{homeomorphisms} $f, g: M \to N$ are said to be \hldef{isotopic} if there exists an isotopy $H: M \times [0, 1] \to N$ such that $h_0 = f$ and $h_1 = g$.

In practice, this terminology is used essentially exclusively in the case that $M$ and $N$ are \CrefAndHyperref{definition:topological_manifold}{topological manifolds with or without boundary}.

The relation of being isotopic is an equivalence relation on the set of homeomorphisms from $M$ to $N$. If $M = N$, the set of isotopy classes of homeomorphisms form a group under composition if the inverse of an isotopy is also an isotopy.
\end{definition}
\begin{definition} \label{definition:isotopy_between_topological_spaces_between_homeomorphisms_relative_to_a_subspace}
    \TODO{do more generally where isotopy might permute some points}
Let $M$ and $N$ be topological spaces, and let $A \subset M$ be a subset. An \hldef{isotopy relative to $A$} is an \CrefAndHyperref{definition:isotopy_between_topological_spaces}{isotopy} $H: M \times [0, 1] \to N$ such that for all $a \in A$ and all $t \in [0, 1]$, $H(a, t) = H(a, 0)$. Two \CrefAndHyperrefIfExist{definition:homeomorphism_of_topological_spaces}{homeomorphisms} $f, g: M \to N$ are \hldef{isotopic relative to $A$} if there exists such an isotopy with $h_0 = f$ and $h_1 = g$.

In practice, this terminology is used essentially exclusively in the case that $M$ and $N$ are \CrefAndHyperref{definition:topological_manifold}{topological manifolds with or without boundary} and $A$ is a submanifold of $M$.
\end{definition}
\begin{definition} \label{definition:smooth_isotopy_on_a_smooth_manifold_between_diffeomorphisms_on_a_smooth_manifold}
Let $M$ be a \CrefAndHyperrefIfExist{definition:C_k_manifold}{smooth manifold}, which may be \CrefAndHyperrefIfExist{definition:topological_manifold}{with or without boundary}. A \hldef{smooth isotopy} is a \CrefAndHyperrefIfExist{definition:C_k_morphism_between_C_k_manifolds}{smooth map} $H: M \times [0, 1] \to M$ such that for each $t \in [0, 1]$, the map $h_t: M \to M$ is a \CrefAndHyperrefIfExist{definition:C_k_morphism_between_C_k_manifolds}{diffeomorphism}. Two diffeomorphisms $f, g: M \to M$ are \hldef{smoothly isotopic} if such a smooth isotopy exists with $h_0 = f$ and $h_1 = g$.

The relation of being isotopic defines an equivalence relation on the group of homeomorphisms $\text{Homeo}(M)$. The set of equivalence classes of $\text{Homeo}(M)$ under this relation is the set of path-components $\pi_0(\text{Homeo}(M))$.
\end{definition}

\subsubsection{Mapping class groups of manifolds}

\begin{definition} \label{definition:compact_open_topology_on_the_set_of_continuous_maps_between_topological_spaces_and_function_space}
    Let $X$ and $Y$ be \CrefAndHyperrefIfExist{definition:topological_space}{topological spaces}. 
    \begin{enumerate}
        \item Write $\mathrm{Open}(Y)$ for the \CrefAndHyperrefIfExist{definition:topological_space}{topology} of $Y$ and $\mathrm{Comp}(X)$ for the family of \CrefAndHyperrefIfExist{compact}{compact subsets} of $X$. 
        For $K \in \mathrm{Comp}(X)$ and $U \in \mathrm{Open}(Y)$, define the subset
        $$\hlin{\langle K,U\rangle \coloneqq \{\, f \in C(X,Y) \mid f(K) \subseteq U \,\} \subseteq C(X,Y)}.$$
        where \CrefAndHyperrefIfExist{definition:continuous_map_of_topological_spaces}{$C(X,Y)$} is the set of all \CrefAndHyperrefIfExist{definition:continuous_map_of_topological_spaces}{continuous maps} $X \to Y$. 
        The collection of all such subsets is denoted by \hl{$\mathcal{S}_{\mathrm{co}}(X,Y)$}.

        \item The \hldef{compact-open topology on $C(X,Y)$} \CrefIfExists{definition:continuous_map_of_topological_spaces} is the \CrefAndHyperrefIfExist{definition:topological_space}{topology} \CrefAndHyperrefIfExist{definition:topology_on_a_set_generated_by_a_collection_of_subsets_and_subbasis_for_a_topology}{generated by} the \CrefAndHyperrefIfExist{definition:topology_on_a_set_generated_by_a_collection_of_subsets_and_subbasis_for_a_topology}{subbasis} $\mathcal{S}_{\mathrm{co}}(X,Y)$. The resulting topological space is called the \hldef{function space from $X$ to $Y$} and is commonly denoted by notations such as \hl{$\operatorname{Map}(X,Y)$}, \hl{$Y^X$}, \hl{$\operatorname{Fun}(X,Y)$}, or \hl{$\operatorname{F}(X,Y)$}. As a set, it equals $C(X,Y)$.

        \item Let $(X,x_0)$ and $(Y,y_0)$ be \CrefAndHyperrefIfExist{definition:pointed_topological_space}{based topological spaces}. The \hldef{based function space from $(X,x_0)$ to $(Y,y_0)$} is the subspace of $\operatorname{Map}(X,Y)$ consisting of all based maps $f$ with $f(x_0)=y_0$, equivalently $C_*((X,x_0),(Y,y_0))$ endowed with the subspace topology inherited from the compact-open function space $\operatorname{Map}(X,Y)$. It itself is a based topological space whose base point is given by the constant map $X \to Y$ with value $y_0$. Common notations for the based function space include those which include a star or bullet such as \hl{$\operatorname{Map}_*((X,x_0),(Y,y_0))$}, \hl{$\operatorname{Map}_*(X,Y)$}, or \hl{$Y^X_*$} to emphasize that the spaces involved are pointed, or those which omit a star or bullet, such as \hl{$\operatorname{Map}(X,Y)$}, and \hl{$Y^X$}. 
    \end{enumerate}
%  Common notations for this function space are \hl{$\operatorname{Map}(X,Y)$} and \hl{$Y^X$}.
\end{definition}

% \begin{definition} \label{definition:compact_open_topology_on_the_set_of_continuous_maps_between_topological_spaces_and_function_space}
%     Let $X,Y$ be \CrefAndHyperrefIfExist{definition:topological_spaces}{topological spaces}.
%     The \hldef{compact-open topology} on the set $C(X,Y)$ of \CrefAndHyperrefIfExist{definition:continuous_map_of_topological_spaces}{continuous maps} is the \CrefAndHyperrefIfExist{definition:basis_for_a_topology}{topology generated by} the \CrefAndHyperrefIfExist{definition:topology_on_a_set_generated_by_a_collection_of_subsets_and_subbasis_for_a_topology}{sub-basis} 
%     $$\{W(K,U) : K \subseteq X \text{ is compact, } U \subseteq Y \text{ is open } \}$$
%     where 
%     $$W(K,U) = \{f \in C(X,y): f(K) \subseteq U\}.$$
% \end{definition}

\begin{definition} \label{definition:homeomorphism_group_of_a_topological_space}
    Let $X$ be a \CrefAndHyperrefIfExist{definition:category_of_topological_spaces}{topological space}. 
    Its \hldef{homeomorphism group} is the group \hl{$\operatorname{Homeo}(X)$} consisting of (self-)\CrefAndHyperrefIfExist{definition:homeomorphism_of_topological_spaces}{homeomorphisms} $X \to X$.

    If $X$ is \CrefAndHyperrefIfExist{definition:locally_compact_topological_space}{locally compact} and \CrefAndHyperrefIfExist{definition:locally_path_connected_topological_space}{locally path-connected}, then $\operatorname{Homeo}(X)$ is a \CrefAndHyperrefIfExist{definition:topological_group}{topological group} equipped with the \CrefAndHyperrefIfExist{definition:compact_open_topology_on_the_set_of_continuous_maps_between_topological_spaces_and_function_space}{compact-open topology} (inherited from the set $C(X,X)$ of continuous maps $X \to X$).
\end{definition}
\begin{notation} \label{notation:homeomorphism_groups_of_a_topological_space_manifolds}
    Let $M$ be a \CrefAndHyperrefIfExist{definition:category_of_topological_spaces}{topological space}.
    % Let $M$ be a \CrefAndHyperrefIfExist{definition:topological_manifold}{topological manifold with or without boundary}. 
    Let $A, B \subseteq M$ be subsets. 
    
    One may denote by notations roughly resembling \hl{$\operatorname{Homeo}(M,A,B)$} for the group of homeomorphisms $f: M \to M$ of $M$ that fix $A$ pointwise and that fix $B$ setwise, i.e. $f|_A = \id_A$ and $f(B) = B$. 
    
    Moreover, 
    \begin{itemize}
        \item When $M$ is an \CrefAndHyperrefIfExist{definition:orientation_of_a_topological_manifold_with_or_without_boundary_using_singular_homology}{oriented} manifold, one adds a superscript of $+$, e.g. \hl{$\operatorname{Homeo}^+(M,A,B)$}, to refer to the subgroup of orientation preserving homeomorphisms.

        \item One adds a subscript of $0$, e.g. \hl{$\operatorname{Homeo}_0(M,A,B)$}, to refer to the subgroup of homeomorphisms $f: M \to M$ that are \CrefAndHyperrefIfExist{definition:isotopy_between_topological_spaces_between_homeomorphisms_relative_to_a_subspace}{isotopic} to $\id_M$ via some isotopy $H: M \times [0,1] \to M$ relative to $A$ such that $H|_{B \times \{t\}}: B \to M$ fixes $B$ setwise for all $t \in [0,1]$. 
    \end{itemize}

    In practice, the following are common notations that may occur:
    \begin{itemize}
        \item \hl{$\operatorname{Homeo}(M, \partial M)$} often denotes $\operatorname{Homeo}(M, \partial M, \emptyset)$, the group of homeomorphisms that fix the boundary of $M$ pointwise. However, when $M$ is a manifold with boundary, \hl{$\operatorname{Homeo}(M)$} may also be used to denote this very same group, which should not be confused with the group of all homeomorphisms of $M$. 

        \item Letting finitely many points $p_1,\ldots,p_n \in M$,  \hl{$\operatorname{Homeo}(M, \{p_1,\ldots,p_n\})$} may denote, depending on the context, either $\operatorname{Homeo}(M, \{p_1,\ldots,p_n\}, \emptyset)$ or $\operatorname{Homeo}(M, \emptyset, \{p_1,\ldots,p_n\})$, which are the groups of homeomorphisms that fix the finitely many points pointwise and setwise, respectively. In particular, elements of the latter group may permute the finitely many points.
    \end{itemize}

    In the case that $M$ is a \CrefAndHyperrefIfExist{definition:topological_manifold}{topological manifold with or without boundary}, $\operatorname{Homeo}(M)$ is a \CrefAndHyperrefIfExist{definition:topological_group}{topological group} equipped with the \CrefAndHyperrefIfExist{definition:compact_open_topology_on_the_set_of_continuous_maps_between_topological_spaces_and_function_space}{compact open topology} inherited from the set $C(M,M)$ of \CrefAndHyperrefIfExist{definition:continuous_map_of_topological_spaces}{continuous maps} $M \to M$. In this case, whenever $A$ is a closed subset of $M$ and $B$ is a finite subset of $M$, the subgroup $\operatorname{Homeo}(M,A,B)$ is in fact a closed subgroup of $\operatorname{Homeo}(M)$. 
\end{notation}

\begin{definition} \label{definition:mapping_class_group_of_a_connected_oriented_topological_manifold_with_or_without_boundary_relative_to_subsets}
Let $M$ be a \CrefAndHyperrefIfExist{definition:disconnected_and_connected_topological_spaces}{connected}, \CrefAndHyperrefIfExist{definition:orientation_of_a_topological_manifold_with_or_without_boundary_using_singular_homology}{oriented} \CrefAndHyperrefIfExist{definition:topological_manifold}{topological $n$-manifold with or without boundary}. Let $A,B \subseteq M$ be subsets. Let \CrefAndHyperref{notation:homeomorphism_groups_of_a_topological_space_manifolds}{$\text{Homeo}^{+}(M, A, B)$} be the group of orientation-preserving homeomorphisms of $M$ that fix $A$ pointwise and that fix $B$ setwise.

The \hldef{(topological) mapping class group of $M$ (with respect to $A$ and $B$)}, which may be denoted \hl{$\text{Mod}(M, A, B)$} or \hl{$\text{Mod}_{\text{Top}}(M, A, B)$}, is the group of \CrefAndHyperref{definition:isotopic_homeomorphisms_between_topological_spaces}{isotopy classes} of these homeomorphisms, defined by the quotient
$$\text{Mod}(M, A, B) = \text{Homeo}^{+}(M, A, B) / \text{Homeo}_{0}(M, A, B)$$
where $\text{Homeo}_{0}(M, A, B)$ is the normal subgroup consisting of homeomorphisms \CrefAndHyperref{definition:isotopy_between_topological_spaces_between_homeomorphisms_relative_to_a_subspace}{isotopic} to the identity via an isotopy fixing $A$ pointwise throughout the isotopy and fixing $B$ setwise throughout the isotopy.


The \hldef{(topological) mapping class group} of $M$ without reference to explicit subsets $A$ and $B$ usually refers to the mapping class group $\text{Mod}(M, \partial M, \emptyset)$. It may often be denoted as \hl{$\text{Mod}(M)$} or as \hl{$\text{Mod}(M, \partial M)$} for emphasis.

In the case that $A$ is a closed subset and $B$ is a finite subset, recall that $\operatorname{Homeo}^+(M,A,B)$ is a closed subgroup of $\operatorname{Homeo}^+(M)$ with the \CrefAndHyperrefIfExist{definition:compact_open_topology_on_the_set_of_continuous_maps_between_topological_spaces_and_function_space}{compact open topology}. In this case, $\text{Mod}(M,A,B)$ is isomorphic to the group $\pi_0(\operatorname{Homeo}^+(M, A,B))$ of path components of $\operatorname{Homeo}^+(M, A,B)$. This group of path components inherits a group structure from the topological group structure on $\operatorname{Homeo}^+(M, A, B)$. 
\end{definition}

\begin{theorem} \label{theorem:restriction_homomorphism_by_taking_out_finitely_many_points_for_mapping_class_group_of_an_oriented_topological_manifold_is_isomorphism_under_nice_conditions}
    Let $M$ be a \CrefAndHyperrefIfExist{definition:disconnected_and_connected_topological_spaces}{connected}, \CrefAndHyperrefIfExist{definition:orientation_of_a_topological_manifold_with_or_without_boundary_using_singular_homology}{oriented} \CrefAndHyperrefIfExist{definition:topological_manifold}{topological $n$-manifold with or without boundary}. Let $A,B \subseteq M$ be subsets.

    The restriction homomorphism
    $$\operatorname{Mod}(M,A,B) \to \operatorname{Mod}(M \setminus B, A \setminus B, \emptyset), \quad f \mapsto f|_{M \setminus B}$$
    is an isomorphism whenever 
    \begin{itemize}
        \item $B$ consists of finitely many points in the interior of $M$,
        \item $M$ is compact,
        \item $\partial M \subseteq A$. 
    \end{itemize}
\end{theorem}

\begin{theorem} \label{theorem:mapping_class_group_on_a_connected_oriented_manifold_with_respect_to_subsets_maps_to_automorphism_group_of_fundamental_group}
    Let $M$ be a \CrefAndHyperrefIfExist{definition:disconnected_and_connected_topological_spaces}{connected}, \CrefAndHyperrefIfExist{definition:orientation_of_a_topological_manifold_with_or_without_boundary_using_singular_homology}{oriented} \CrefAndHyperrefIfExist{definition:topological_manifold}{topological $n$-manifold with or without boundary}. Let $A,B \subseteq M$ be subsets. Let $x_0 \in A$ be a fixed basepoint of $M$. 

    \begin{enumerate}
        \item 
        There is a well-defined group homomorphism
        $$\Phi: \Mod(M,A,B) \to \operatorname{Aut}(\pi_1(M,x_0))$$ 
        (\Cref{definition:mapping_class_group_of_a_connected_oriented_topological_manifold_with_or_without_boundary_relative_to_subsets})\CrefIfExists{definition:homotopy_groups_of_a_pointed_topological_space}
        defined by $[f] \mapsto f_*$ where $f_*$ is the automorphism $\pi_1(M,x_0) \to \pi_1(M,x_0)$ induced by $f: (M,x_0) \to (M,x_0)$. In particular, isotopic elements of $\operatorname{Homeo}^+(M,A,B)$ (\Cref{notation:homeomorphism_groups_of_a_topological_space_manifolds}) are sent to the same automorphism on $\pi_1(M,x_0)$.
        \TODO{If $x_0$ is not fixed pointwise, then one only gets a map to the outer automorhpism group.}

        \item Assuming that $x_0 \not\in B$, there is a well-defined group homomorphism
        $$\Phi: \Mod(M,A,B) \to \operatorname{Aut}(\pi_1(M \setminus B,x_0))$$ 
        defined by $[f] \mapsto (f|_{M \setminus B})_*$. This follows from the fact that any representative $f$ restricts to a pointed homeomorphism of the complement.
    \end{enumerate}
\end{theorem}
\begin{proof}

    \begin{enumerate}
        \item  \TODO{}
        \item  Take the composition 
        $$\Mod(M,A,B) \xrightarrow{\operatorname{res}} \Mod(M \setminus B, A \setminus B, \emptyset) \xrightarrow{\Phi'} \Aut(\pi_1(M \setminus B, x_0));$$
        where $\operatorname{res}$ is the homomorhpism induced by restricting homeomorphisms and isotopies to the complement, and $\Phi'$ is the homomorphism obtained by applying the previous part to the triple $(M \setminus B, A \setminus B, \emptyset)$.
    \end{enumerate}

\end{proof}


\TODO{}
\begin{notation}
Let \hl{$S_{g,n}^b$} denote a connected, orientable surface (a 2-manifold) of genus $g$ with $n$ punctures and $b$ boundary components. We denote the mapping class group of this surface as \hl{$\text{Mod}(S_{g,n}^b)$}. When $n=0$ and $b=0$, we write \hl{$\text{Mod}(S_g)$}.
\end{notation}

\TODO{}
\begin{definition}
Let $S$ be an orientable surface and let $c$ be a simple closed curve on $S$. A \hldef{Dehn twist about $c$}, denoted \hl{$T_c$}, is an element of the mapping class group $\text{Mod}(S)$ represented by a homeomorphism obtained by cutting $S$ along $c$, rotating one boundary of the cut by $2\pi$, and reglueing.
\end{definition}

\TODO{}
\begin{theorem}
Let $S_{g,n}^b$ be a connected, orientable surface of finite type. For any integers $g, n, b \geq 0$, the mapping class group $\text{Mod}(S_{g,n}^b)$ is finitely generated. Specifically, if $b=0$ and $n=0$, the group $\text{Mod}(S_g)$ is generated by a finite set of Dehn twists about a collection of simple closed curves.
\end{theorem}


\TODO{}
\begin{proposition}
Let $S$ be a connected, orientable surface of genus $g$ with $n$ punctures and $b$ boundary components such that its Euler characteristic $\chi(S)$ is negative. The mapping class group $\text{Mod}(S)$ acts on the Teichmüller space $\mathcal{T}(S)$ by precomposition.
\textit{Remark:} \TODO{define Teichmüller space}.
\end{proposition}

\TODO{}
\begin{theorem}
Let $S$ be a connected, orientable surface of finite type with $\chi(S) < 0$. The mapping class group $\text{Mod}(S)$ acts properly discontinuously on the Teichmüller space $\mathcal{T}(S)$. The quotient space
$$\hlin{\mathcal{M}(S) = \mathcal{T}(S) / \text{Mod}(S)}$$
is the moduli space of Riemann surfaces homeomorphic to $S$.
\textit{Remark:} \TODO{define moduli space of Riemann surfaces}.
\end{theorem}


\begin{definition} \label{definition:smooth_mapping_class_group_of_a_smooth_orientable_manifold_with_or_without_boundary}
Let $M$ be a smooth orientable $n$-manifold, which may be with or without boundary. The \hldef{smooth mapping class group of $M$}, denoted \hl{$\text{Mod}(M)$} or \hl{$\text{Mod}_{\text{smooth}}(M)$}, not to be confused with the \CrefAndHyperrefIfExist{definition:mapping_class_group_of_a_connected_oriented_topological_manifold_with_or_without_boundary_relative_to_subsets}{topological mapping class group of $M$}, is the group of \CrefAndHyperrefIfExist{definition:isotopy_between_topological_spaces_between_homeomorphisms_relative_to_a_subspace}{isotopy classes} of \CrefAndHyperrefIfExist{definition:orientation_preserving_homeomorphism_group_fixing_boundary_of_an_oriented_topological_manifold_with_or_without_boundary}{orientation-preserving} \CrefAndHyperrefIfExist{definition:C_k_morphism_between_C_k_manifolds}{diffeomorphisms} of $M$ that fix the \CrefAndHyperrefIfExist{definition:interior_and_boundary_of_topological_space}{boundary} $\partial M$ pointwise:
$$\text{Mod}(M) = \pi_0(\text{Diff}^{+}(M, \partial M))$$
\CrefIfExists{definition:path_component_of_a_point_in_a_topological_space}

\TODO{separate notation}
where $\text{Diff}^{+}(M, \partial M)$ is the topological group of orientation-preserving diffeomorphisms restricting to the identity on $\partial M$, equipped with the compact-open topology.
\end{definition}

\TODO{laruent polynomial ring}


\subsection{The Burau representations}



\begin{definition} \label{definition:representation_of_a_mapping_class_group_on_the_singular_homology_groups_of_a_cover}
    \TODO{At the moment, the ``reduced'' version is being described. The ``unreduced'' version needs to be described, which should be based upon havnig $H_n(\tilde{M}, p^{-1}(x_0)$. The Hurewicz morhpism should also be discussed}
    Let $M$ be a \CrefAndHyperrefIfExist{definition:disconnected_and_connected_topological_spaces}{connected}, \CrefAndHyperrefIfExist{definition:orientation_of_a_topological_manifold_with_or_without_boundary_using_singular_homology}{oriented} \CrefAndHyperrefIfExist{definition:topological_manifold}{topological $n$-manifold with or without boundary}. Let $A,B \subseteq M$ be subsets. Let $x_0 \in A$ be a fixed basepoint of $M$. Let $p: (\tilde{M}, \tilde{x}_0) \to (M,x_0)$ be a \CrefAndHyperrefIfExist{definition:covering_space_of_a_topological_space}{pointed covering}. Write $H$ for the subgroup of $\pi_1(M,x_0)$ \CrefAndHyperrefIfExist{proposition:isom_class_of_conn_pointed_covers_of_a_conn_loc_path_conn_semiloc_simp_conn_pointed_space_corresponds_to_subgroup_of_fund_gp}{corresponding to} the covering $p$.
    
    Note that there is a ``forgetful'' group homomorphism 
    $$\operatorname{Homeo}^+(M,A,B) \to \Aut(M,x_0)$$
    (\Cref{notation:homeomorphism_groups_of_a_topological_space_manifolds}). 
    Assume for all $f \in \operatorname{Homeo}^+(M,A,B)$ (\Cref{notation:homeomorphism_groups_of_a_topological_space_manifolds}) the following ``lifting criterion'' holds: $f_*(H) \subseteq H$. 

    As a consequence of \Cref{theorem:lift_for_a_map_of_pointed_spaces_to_a_map_of_covers_exists_iff_images_of_fundamental_groups_are_contained}, there is then a group homomorphism
    $$\operatorname{Homeo}^+(M,A,B) \to \Aut_{\Top_*}(\tilde{M},\tilde{x}_0).$$
    \CrefIfExists{definition:pointed_topological_space}

    For $n \geq 1$, 
    One can show that the composition 
    $$\operatorname{Homeo}^+(M,A,B) \to \Aut_{\Top_*}(\tilde{M},\tilde{x}_0) \to \Aut(\pi_n(\tilde{M}, \tilde{x}_0))$$
    sends \CrefAndHyperrefIfExist{definition:isotopy_between_topological_spaces_between_homeomorphisms_relative_to_a_subspace}{isotopic} homeomorphisms to the same elements of $\Aut(\pi_n(\tilde{M}, \tilde{x}_0))$ and hence the composition factors through $\operatorname{Mod}(M,A,B)$:
    \TODO{Okay, so to get the action on $H_n$, I really should not go through $\pi_n$ because an action on $\pi_n$ is not generally sufficient to get an action on $H_n$. However, since $H_n$ is functorial, I can just directly get the action on $H_n$ to start with, and argue that $H_n$ is ``homotopy invariant'', thereby obtaining a well defined map on $\operatorname{Mod}(M,A,B)$}
    $$\operatorname{Mod}(M,A,B) \to \Aut(\pi_n(\tilde{M}, \tilde{x}_0)).$$
    By the naturality of the Hurewicz homomorphism \TODO{} $\pi_n(\tilde{M}, \tilde{x}_0) \to H_n(\tilde{M}, \tilde{x}_0; \bbZ)$ and the functoriality of singular homology, this action induces a well-defined representation
    $$\operatorname{Mod}(M,A,B) \to \operatorname{Aut}(H_n(\tilde{M}, \tilde{x}_0; \bbZ)).$$
    on $H_n(\tilde{M}, \tilde{x}_0; \bbZ)$.

    In fact, for general abelian groups $G$, these representations may be extended to representations
    $$\operatorname{Mod}(M,A,B) \to \operatorname{Aut}(H_n(\tilde{M}, \tilde{x}_0; G))$$
    as follows: by the \CrefAndHyperrefIfExist{theorem:universal_coefficient_theorem_for_singular_homology_of_topological_spaces_with_coefficients_in_abelian_groups}{universal coefficient theorem}, there are short exact sequences
    $$0 \to H_n(\tilde{M}, \tilde{x}_0; \bbZ) \otimes G \to H_n(\tilde{M}, \tilde{x}_0; G) \to \Tor_1^{\bbZ}(H_{n-1}(\tilde{M}, \tilde{x}_0;\bbZ), G) \to 0$$
    natural in $(\tilde{M}, \tilde{x}_0)$. We have the group homomorphism
    $$\operatorname{Homeo}^+(M,A,B) \to \Aut_{\Top_*}(\tilde{M},\tilde{x}_0) \to \Aut(H_n(\tilde{M}, \tilde{x}_0; G)).$$
    To see that this group homomorphism sends isotopic elements of $\operatorname{Homeo}^+(M,A,B)$ to the same automorphism on $H_n(\tilde{M}, \tilde{x}_0; G)$, the five lemma may be used on the short exact sequence.

    The following names seem appropriate for the above representation:
    \begin{enumerate}
        \item \hldef{The representation of $\operatorname{Mod}(M,A,B)$ on the $n$th singular homology of the cover $(\tilde{M}, \tilde{x}_0)$}
        \item A \hldef{generalized Lawrence-Krammer-Bigelow type representation}
    \end{enumerate}
    If $G = R$ is a ring, then the representations are representations over $R$.
    \TODO{representation}
    % Let $\Top_{*,\alpha}$ be the category of pointed spaces $(X,x_0)$ equipped with a surjective homomorphism $\alpha: \pi_1(X,x_0) \to \bbZ$. There is a functor $\Top_{*,\alpha}$ 
\end{definition}


\begin{definition}
    Let $\bbD \subseteq \bbR^2$ be a (closed unit) disk, let $x_1,\ldots,x_n \in \bbD$ be $n$-distinct points. Let $G$ be an abelian group (or ring). 
    The \hldef{unreduced Burau representation}
\end{definition}

\TODO{define burau, gassner}

\TODO{unreduced, reduced Burau representation}


\TODO{relationship between $H_1(D)$, $H_1(D, x_0)$, $H_1(\tilde{D}, \tilde{x}_0)$, $H_1(\tilde{D}, p^{-1}(\tilde{x}_0))$}


\begin{definition}[Homotopy groups of a pointed topological space] \label{definition:homotopy_groups_of_a_pointed_topological_space}
For any \CrefAndHyperrefIfExist{definition:pointed_topological_space}{pointed topological space} $(X,x_0)$ and integer $n \ge 0$, the \hldef{$n$-th homotopy set of $X$ at $x_0$}, denoted \hl{$\pi_n(X,x_0)$} or simply \hl{$\pi_n(X)$} when the base point $x_0$ is understood or a choice of base point does not really matter, is defined as the set of all \CrefAndHyperrefIfExist{definition:homotopy_class_of_maps_of_topological_spaces_relative_to_a_subset}{homotopy classes (rel.\ $\partial I^n$)} of \CrefAndHyperrefIfExist{definition:map_of_pairs_of_topological_spaces}{paired maps}
$$f : (I^n, \partial I^n) \to (X,\{x_0\}),$$
where $I^n = [0,1]^n$ and $\partial I^0$ is taken to be the singleton $I^0$ instead of the literal \CrefAndHyperrefIfExist{definition:interior_and_boundary_of_topological_space}{boundary} thereof. Equivalently, one can also define it as the set of all homotopy classes of \CrefAndHyperrefIfExist{definition:pointed_topological_space}{pointed maps}
$$f: (S^n, s_0) \to (X,x_0),$$
where $S^n$ is the $n$-sphere \TODO{} and $s_0 \in S^n$ is a choice of base point. In other words, $\pi_n(X,x_0)$ is by definition $[(S^n, s_0), (X,x_0)]_{*}$\CrefIfExists{definition:homotopy_class_of_maps_of_topological_spaces_relative_to_a_subset}.

For $n \ge 1$, $\pi_n(X,x_0)$ is a group under concatenation of pointed maps, and for $n \ge 2$, it is abelian. As such, for $n \geq 1$, it is appropriate to refer to $\pi_n(X,x_0)$ as the \hldef{$n$th homotopy group of $(X,x_0)$}.

The \hldef{fundamental group of $(X,x_0)$} refers to $\pi_1(X,x_0)$. Equivalently, it is the group of homotopy classes (rel.\ endpoints) of \CrefAndHyperrefIfExist{definition:path_and_loop_in_a_topological_space}{loops} $\gamma : [0,1] \to X$ satisfying $\gamma(0)=\gamma(1)=x_0$. 

For every $n \ge 0$, the assignment $(X,x_0) \mapsto \pi_n(X,x_0)$ is functorial --- in fact
\begin{enumerate}
    \item For $n = 0$, it is a functor $\mathrm{Top}_* \to \Sets$\CrefIfExists{definition:category_of_sets},
    \item For $n = 1$, it is a functor $\mathrm{Top}_* \to \mathbf{Grp}$\CrefIfExists{definition:category_of_groups_of_abelian_groups},
    \item For $n \geq 2$, it is a functor $\mathrm{Top}_* \to \mathbf{Ab}$(\CrefIfExists{definition:category_of_groups_of_abelian_groups}).
\end{enumerate}

Given a map $f: (X,x_0) \to (Y,y_0)$ of pointed spaces, the induced map $\pi_n(X,x_0) \to \pi_n(Y,y_0)$ is often denoted by \hl{$f_*$}.

\TextIfExists{definition:relative_homotopy_groups_of_pointed_pairs_of_topological_spaces}{
The homotopy set $\pi_n(X,x_0)$ is naturally isomorphic to the \CrefAndHyperrefIfExist{definition:relative_homotopy_groups_of_pointed_pairs_of_topological_spaces}{relative homotopy set $\pi_n(X,\{x_0\},x_0)$}.
}

We further note that the set $\pi_n(X,x_0) = [(S^n, s_0), (X,x_0)]_{*}$ is identifiable with $\pi_0(\operatorname{Map}_*((S^n, s_0), (X,x_0)))$\CrefIfExists{proposition:based_homotopy_classes_of_maps_between_based_spaces_correspond_to_path_components_of_function_space}. 

\end{definition}


\section{Hurwitz spaces}


\subsection{Topological Hurwitz space}


\subsection{Algebro-geometric Hurwitz space}


\subsubsection{Nielsen classes}

\subsubsection{Racks and quandles}

\begin{definition} \label{definition:rack} 	
    A \hldef{rack} is a set $c$ equipped with an \hldef{action map} 
    $$\hlin{\triangleright: c \times c 	\to c, \quad (a,b) \mapsto a \triangleright b}$$
    satisfying the following equivalent definitions:
    \begin{enumerate}
        \item for all $n \geq 1$ and all $1 \leq i \leq n-1$, the operation
        \begin{align*} 		\sigma_i : c^n & \rightarrow c^n \\ 		(x_1, \ldots, x_{i-1}, x_i, x_{i+1}, x_{i+2}, \ldots, x_n) & \mapsto 		(x_1, \ldots, x_{i-1}, x_{i+1}, x_{i+1} \triangleright x_i, x_{i+2}, \ldots, 		x_n) 	\end{align*}
        defines an action of the \CrefAndHyperrefIfExist{definition:artin_braid_group}{braid group $B_n$} where $\sigma_i \in B_n$ is the $i$th elementary braid.
        
        \item the map $x \triangleright (-): c \to c$ is a bijection for each $x \in c$ and $x \triangleright (y \triangleright z) = (x \triangleright y) \triangleright (x \triangleright z)$. 
    \end{enumerate}
\end{definition}
\begin{definition} \label{definition:homomorphism_of_racks}
Let $(X, \rhd)$ and $(Y, \rhd')$ be \CrefAndHyperrefIfExist{definition:rack}{racks}. 
A \hldef{homomorphism of racks} is a map $f: X \to Y$ such that for all $x_1, x_2 \in X$, the condition $f(x_1 \rhd x_2) = f(x_1) \rhd' f(x_2)$ holds.
\end{definition}
\begin{definition} \label{definition:category_of_racks}
    The \hldef{category of racks} has as objects \CrefAndHyperrefIfExist{definition:rack}{racks} and as morphisms the \CrefAndHyperrefIfExist{definition:homomorphism_of_racks}{homomorphisms between racks}. This category may be denoted \hl{$\mathbf{Rack}$}.
\end{definition}

\begin{definition} \label{definition:left_translation_map_of_a_rack_by_an_element}
Let $(X, \rhd)$ be a \CrefAndHyperrefIfExist{definition:rack}{rack}. For each $x \in X$, the \hldef{left translation map} \hl{$L_x: X \to X$}, defined by $L_x(y) = x \rhd y$, is an \CrefAndHyperrefIfExist{definition:homomorphism_of_racks}{automorphism of the rack}. 
If $(X, \rhd)$ is a \CrefAndHyperrefIfExist{definition:quandle}{quandle}, then the idempotency axiom implies that $x$ is a fixed point of $L_x$ for all $x \in X$.
\end{definition}
\begin{definition} \label{definition:rack_of_a_conjugacy_invariant_subset_of_a_group}
Let $G$ be a \CrefAndHyperrefIfExist{definition:group}{group} and $S \subset G$ a \CrefAndHyperrefIfExist{definition:subset_of_a_set}{subset} that is \CrefAndHyperrefIfExist{definition:conjugacy_invariant_subset_of_a_group}{invariant under conjugation}, such that for all $g \in G$ and $s \in S$, $g s g^{-1} \in S$. 
The \hldef{rack of a conjugacy invariant subset} is the \CrefAndHyperrefIfExist{definition:rack}{rack} with underlying set $S$ equipped with the operation $x \rhd y = x y x^{-1}$ for all $x, y \in S$.

In the case that $S = G$, one might call this rack the \hldef{conjugation rack} or \hldef{conjugation quandle} (\Cref{definition:quandle}) of $G$ and denote by \hl{$\operatorname{Conj}(G)$}. Note that the assignment of a \CrefAndHyperrefIfExist{definition:group}{group} to its \CrefAndHyperrefIfExist{definition:conjugation_of_group_elements_and_subgroup_by_a_group_element}{conjugation} rack is \CrefAndHyperrefIfExist{definition:functor_between_categories}{functorial} from the \CrefAndHyperrefIfExist{definition:category_of_groups_of_abelian_groups}{category of groups} to the \CrefAndHyperrefIfExist{definition:category_of_racks}{category of racks}.
\end{definition}
\begin{definition} \label{definition:automorphism_group_of_a_rack}
Let $(X, \rhd)$ be a \CrefAndHyperrefIfExist{definition:rack}{rack}. 
The \hldef{automorphism group of the rack}, denoted \hl{$\mathrm{Aut}(X)$}, is the set of all \CrefAndHyperrefIfExist{definition:injective_surjective_bijective_map_of_sets}{bijections} $f: X \to X$ such that $f(x \rhd y) = f(x) \rhd f(y)$ for all $x, y \in X$, forming a \CrefAndHyperrefIfExist{definition:group}{group} under composition.
\end{definition}
\begin{definition} \label{definition:inner_automorphism_group_of_a_rack}
Let $(X, \rhd)$ be a \CrefAndHyperrefIfExist{definition:rack}{rack}. 
The \hldef{inner automorphism group} \hl{$\mathrm{Inn}(X)$} is the \CrefAndHyperrefIfExist{definition:subgroup_of_a_group}{subgroup} of the \CrefAndHyperrefIfExist{definition:symmetric_group_on_a_set}{permutation group} $\mathrm{Sym}(X)$ generated by the \CrefAndHyperrefIfExist{definition:left_translation_map_of_a_rack_by_an_element}{left translations} $L_x: X \to X$ defined by $L_x(y) = x \rhd y$ for all $x, y \in X$.
\end{definition}
\begin{definition} \label{definition:reduced_structure_group_of_a_rack}
Let $(X, \rhd)$ be a \CrefAndHyperrefIfExist{definition:rack}{rack} and let $F(X)$ be the \CrefAndHyperrefIfExist{definition:free_group_generated_by_a_set}{free group generated by} the set $X$. 
The \hldef{reduced structure group of a rack}, often denoted \hl{$G_X$}, is the \CrefAndHyperrefIfExist{definition:quotient_of_a_group_by_a_normal_subgroup}{quotient} of $F(X)$ by the normal subgroup generated by the relations $x \rhd y = x y x^{-1}$ for all $x, y \in X$.
\TODO{normal subgroup generated by}

The reduced structure group may also be called the \hldef{adjoint group} or \hldef{structure group}.
\end{definition}
\begin{proposition} \label{proposition:reduced_structure_group_of_a_rack_is_left_adjoint_to_conjugation_rack_of_a_group}
Let $(X, \rhd)$ be a \CrefAndHyperrefIfExist{definition:rack}{rack} and let $G_X$ be its \CrefAndHyperrefIfExist{definition:reduced_structure_group_of_a_rack}{reduced structure group}.
There exists a unique \CrefAndHyperrefIfExist{definition:homomorphism_of_racks}{rack homomorphism} $\phi: X \to G_X$ where $G_X$ is \CrefAndHyperrefIfExist{definition:rack_of_a_conjugacy_invariant_subset_of_a_group}{viewed as a rack} under \CrefAndHyperrefIfExist{definition:conjugation_of_group_elements_and_subgroup_by_a_group_element}{conjugation}, such that for any group $H$ and any rack homomorphism $f: X \to H$, there exists a unique \CrefAndHyperrefIfExist{definition:group_homomorphism}{group homomorphism} $\bar{f}: G_X \to H$ satisfying $f = \bar{f} \circ \phi$.

In other words, the \CrefAndHyperrefIfExist{definition:functor_between_categories}{functor} from the \CrefAndHyperrefIfExist{definition:category_of_racks}{category of racks} to the \CrefAndHyperrefIfExist{definition:category_of_groups_of_abelian_groups}{category of groups} sending a rack $X$ to $G_X$ is \CrefAndHyperrefIfExist{definition:adjoint_functors_between_categories_unit_counit_of_adjoint_functors}{left adjoint} to the functor sending a group to its \CrefAndHyperrefIfExist{definition:rack_of_a_conjugacy_invariant_subset_of_a_group}{conjugacy rack}.
\end{proposition}



\begin{definition} \label{definition:quandle}
A \hldef{quandle} is a \CrefAndHyperrefIfExist{definition:rack}{rack} $(X, \rhd)$ that satisfies the idempotency axiom: for all $x \in X$, $x \rhd x = x$.
\end{definition}
\begin{definition} \label{definition:homomorphism_of_quandles}
Let $(X, \rhd)$ and $(Y, \rhd')$ be \CrefAndHyperrefIfExist{definition:quandle}{quandles}. A \hldef{homomorphism of quandles} is a map $f: X \to Y$ that is a \CrefAndHyperrefIfExist{definition:homomorphism_of_racks}{homomorphism of the underlying racks}, i.e. a map satisfying $f(x_1 \rhd x_2) = f(x_1) \rhd' f(x_2)$ for all $x_1, x_2 \in X$.
\end{definition}
\begin{definition} \label{definition:category_of_quandles}
The \hldef{category of quandles}, denoted \hl{$\mathbf{Qnd}$}, is the \CrefAndHyperrefIfExist{definition:full_subcategory_of_a_category}{full subcategory} of the \CrefAndHyperrefIfExist{definition:category_of_racks}{category of racks} whose objects are the \CrefAndHyperrefIfExist{definition:quandle}{quandles}. 
\end{definition}
\begin{definition} \label{definition:quandleization_of_a_rack}
    \TODO{quotient of a set by equivalence relation}
    \TODO{equivalence generated by relations}
    Given a \CrefAndHyperrefIfExist{definition:rack}{rack} $(X, \rhd)$, its \hldef{quandle-ization} is the \CrefAndHyperrefIfExist{definition:quandle}{quandle} on $X/\sim$ whose \CrefAndHyperrefIfExist{definition:rack}{action map} is inherited from $\rhd$, where $\sim$ is the congruence on $X$ generated by the relations $X \rhd x \sim x$ for all $x \in X$. The assignment sending a rack to its quandle-ization is \CrefAndHyperrefIfExist{definition:functor_between_categories}{functorial}. 
\end{definition}
\begin{proposition} \label{proposition:quandleization_fucntor_is_left_adjoint_to_inclusion_of_category_of_quandles_in_racks}
    The inclusion functor $\iota: \mathbf{Qnd} \hookrightarrow \mathbf{\CrefAndHyperrefIfExist{definition:rack}{Rack}}$ is  \CrefAndHyperrefIfExist{definition:adjoint_functors_between_categories_unit_counit_of_adjoint_functors}{right adjoint} to the \CrefAndHyperrefIfExist{definition:quandleization_of_a_rack}{quandle-ization functor} $Q: \mathbf{Rack} \to \mathbf{Qnd}$. 
\end{proposition}


\begin{definition} \label{definition:set_of_components_of_a_rack}
    Let $c$ be a \CrefAndHyperrefIfExist{definition:rack}{rack}. 
    There is a rack whose underlying set is the \CrefAndHyperrefIfExist{definition:orbit_of_a_group_action}{orbit} of $c$ under the action of the \CrefAndHyperrefIfExist{definition:reduced_structure_group_of_a_rack}{reduced structure group} of $c$ and whose action map is \CrefAndHyperrefIfExist{definition:topos_morphism_induced_by_a_morphism_of_small_sites}{given by} $x \rhd y = y$ for all $x,y \in c/c$. We may refer to this rack $c/c$ as the \hldef{set of components of the rack $c$}, the \hldef{set of orbits of the rack $c$}, the \hldef{orbit rack of $c$}, or the \hldef{constant rack of $c$}.
\end{definition}



\section{Lifting invariants}




\section{Homological stability of Hurwitz spaces}









\appendix

\section{Connectivity of topological spaces}

\begin{definition} \label{definition:path_connected_topological_space}
    Let $X$ be a \CrefAndHyperrefIfExist{definition:category_of_topological_spaces}{topological space}.  
    It is said to be \hldef{path connected} if for any two points $x_0, x_1$ in $X$, there is some \CrefAndHyperrefIfExist{definition:path_and_loop_in_a_topological_space}{path} $\gamma: [0,1] \to X$ in $X$ such that $\gamma(0) = x_0$ and $\gamma(1) = x_1$. 
\end{definition}
\begin{definition} \label{definition:locally_path_connected_topological_space}
    Let $X$ be a \CrefAndHyperrefIfExist{definition:category_of_topological_spaces}{topological space}.  
    It is said to be \hldef{locally path connected} if for any $x \in X$, there is an \CrefAndHyperrefIfExist{definition:neighborhood_and_neighborhood_basis_of_a_point_in_a_topological_space}{open neighborhood} $x \in U \subseteq X$ that is \CrefAndHyperrefIfExist{definition:path_connected_topological_space}{path connected}.
\end{definition}
\begin{definition} \label{definition:path_component_of_a_point_in_a_topological_space}
    Let $X$ be a \CrefAndHyperrefIfExist{definition:category_of_topological_spaces}{topological space}. 
    
    Let $x \in X$. The \hldef{path component of $x$ in $X$} is the maximal path connected subspace of $X$ containing $x$.

    The set of path components in $X$ is denoted \hl{$\pi_0(X)$}. For any $x \in X$, $\pi_0(X)$ coincides, as a set, with the \CrefAndHyperrefIfExist{definition:homotopy_groups_of_a_pointed_topological_space}{zeroth homotopy set} $\pi_0(X,x_0)$ of the \CrefAndHyperrefIfExist{definition:pointed_topological_space}{pointed space} $(X,x_0)$.
\end{definition}
\begin{definition} \label{definition:simply_connected_space}
A \CrefAndHyperrefIfExist{definition:category_of_topological_spaces}{topological space} $X$ is \hldef{simply connected} if it is \CrefAndHyperrefIfExist{definition:path_connected_topological_space}{path-connected} and its \CrefAndHyperrefIfExist{definition:homotopy_groups_of_a_pointed_topological_space}{fundamental group} $\pi_1(X, x_0)$ is trivial for some (and hence every) \CrefAndHyperrefIfExist{definition:basis_for_a_topology}{base} \CrefAndHyperrefIfExist{definition:point_of_a_topos}{point} $x_0 \in X$. Equivalently, $X$ is simply connected if any two \CrefAndHyperrefIfExist{definition:continuous_map_of_topological_spaces}{continuous maps} $f, g: [0, 1] \to X$ with $f(0) = g(0)$ and $f(1) = g(1)$ are \CrefAndHyperrefIfExist{definition:homotopy_of_maps_of_topological_spaces_relative_to_a_subset}{homotopic relative to} $\{0, 1\}$.
\end{definition}
\begin{proposition} \label{proposition:topological_space_is_simply_connected_iff_path_connected_and_every_continuous_map_from_the_circle_is_null_homotopic}
Let $X$ be a \CrefAndHyperrefIfExist{definition:category_of_topological_spaces}{topological space}. $X$ is \CrefAndHyperrefIfExist{definition:simply_connected_space}{simply connected} if and only if it is \CrefAndHyperrefIfExist{definition:path_connected_topological_space}{path-connected} and every \CrefAndHyperrefIfExist{definition:continuous_map_of_topological_spaces}{continuous map} $S^1 \to X$ from the unit circle into $X$ is null-homotopic.
\end{proposition}

\begin{definition} \label{definition:locally_simply_connected_topological_space}
A \CrefAndHyperrefIfExist{definition:category_of_topological_spaces}{topological space} $X$ is \hldef{locally simply connected} if it admits a \CrefAndHyperrefIfExist{definition:basis_for_a_topology}{basis} of \CrefAndHyperrefIfExist{definition:topological_space}{open sets} $\{U_{\alpha}\}$ such that each $U_{\alpha}$ is simply \CrefAndHyperrefIfExist{definition:disconnected_and_connected_topological_spaces}{connected}. That is, for every $x \in X$ and every \CrefAndHyperrefIfExist{definition:neighborhood_and_neighborhood_basis_of_a_point_in_a_topological_space}{open neighborhood} $V$ of $x$, there exists a simply connected open neighborhood $U \subset V$ of $x$.
\end{definition}
\begin{proposition} \label{proposition:locally_simply_connected_topological_space_is_locally_path_connected_and_semi_locally_simply_connected}
Let $X$ be a \CrefAndHyperrefIfExist{definition:category_of_topological_spaces}{topological space}. 
If $X$ is \CrefAndHyperrefIfExist{definition:locally_simply_connected_topological_space}{locally simply connected}, then $X$ is both \CrefAndHyperrefIfExist{definition:locally_path_connected_topological_space}{locally path-connected} and \CrefAndHyperrefIfExist{definition:semi_locally_simply_connected_space}{semi-locally simply connected}.
\end{proposition}
\begin{proposition} \label{proposition:topological_manifolds_are_locally_simply_connected}
Every \CrefAndHyperrefIfExist{definition:topological_manifold}{topological manifold, with or without boundary}, is \CrefAndHyperrefIfExist{definition:locally_simply_connected_topological_space}{locally simply connected} because every point has a neighborhood homeomorphic to an open ball in $\mathbb{R}^n$ or $\mathbb{H}^n$, both of which are simply connected. 
\end{proposition}
\begin{definition} \label{definition:locally_contractible_topological_space}
    \TODO{contractible}
Let $X$ be a \CrefAndHyperrefIfExist{definition:category_of_topological_spaces}{topological space}. $X$ is \hldef{locally contractible} if every point $x \in X$ has a \CrefAndHyperrefIfExist{definition:basis_for_a_topology}{basis} of contractible \CrefAndHyperrefIfExist{definition:neighborhood_and_neighborhood_basis_of_a_point_in_a_topological_space}{open neighborhoods}. Since every contractible space is simply connected, every locally contractible space is \CrefAndHyperrefIfExist{definition:locally_simply_connected_topological_space}{locally simply connected}.
\TODO{simply connected space}
\end{definition}
\begin{definition} \label{definition:semi_locally_simply_connected_space}
A \CrefAndHyperrefIfExist{definition:category_of_topological_spaces}{topological space} $X$ is \hldef{semi-locally simply connected} if for every point $x \in X$, there exists an \CrefAndHyperrefIfExist{definition:neighborhood_and_neighborhood_basis_of_a_point_in_a_topological_space}{open neighborhood} $U$ of $x$ such that the inclusion-induced \CrefAndHyperrefIfExist{definition:group_homomorphism}{homomorphism} on \CrefAndHyperrefIfExist{definition:homotopy_groups_of_a_pointed_topological_space}{fundamental groups}
$$i_*: \pi_1(U, x) \to \pi_1(X, x)$$
is the trivial homomorphism.
\end{definition}
\begin{proposition} \label{proposition:topological_space_is_semi_local_simply_connected_iff_every_point_has_open_neighborhood_such_that_every_loop_based_there_is_null_homotopic_in_ambient_space}
Let $X$ be a \CrefAndHyperrefIfExist{definition:category_of_topological_spaces}{topological space}. 
The condition of \CrefAndHyperrefIfExist{definition:semi_locally_simply_connected_space}{semi-local simple connectivity} is equivalent to the requirement that every point $x \in X$ has an open neighborhood $U$ such that every loop in $U$ based at $x$ is null-homotopic in $X$. Note that the loops are not required to be null-homotopic within $U$ itself.
\TODO{null-homotopic}
\end{proposition}

\section{Covering spaces}

\begin{definition} \label{definition:covering_space_of_a_topological_space}
    \begin{enumerate}
        \item Let $X$ be a \CrefAndHyperrefIfExist{definition:category_of_topological_spaces}{topological space}. 

        A \hldef{covering space of $X$} is a pair $(\tilde{X}, p)$ where $\tilde{X}$ is a topological space and $p: \tilde{X} \to X$ is a \CrefAndHyperrefIfExist{definition:continuous_map_of_topological_spaces}{continuous} \CrefAndHyperrefIfExist{definition:injective_surjective_bijective_map_of_sets}{surjective} map such that for every $x \in X$, there exists an \CrefAndHyperrefIfExist{definition:neighborhood_and_neighborhood_basis_of_a_point_in_a_topological_space}{open neighborhood} $U$ of $x$ (called an \hldef{evenly covered neighborhood)} for which the preimage $p^{-1}(U)$ is a \CrefAndHyperrefIfExist{definition:disjoint_union_of_topological_spaces}{disjoint union} of \CrefAndHyperrefIfExist{definition:topological_space}{open sets} $V_i \subset \tilde{X}$, each of which is mapped \CrefAndHyperrefIfExist{definition:homeomorphism_of_topological_spaces}{homeomorphically} onto $U$ by the restriction $p|_{V_i}$.

        \item Let $(X,x_0)$ be a \CrefAndHyperrefIfExist{definition:pointed_topological_space}{pointed topological space}. A \hldef{(pointed) covering of $(X,x_0)$} is a \CrefAndHyperrefIfExist{definition:pointed_topological_space}{pointed topological space} $(\tilde{X}, \tilde{x}_0)$ along with a covering $p: \tilde{X} \to X$ of $X$ such that $p$ is a \CrefAndHyperrefIfExist{definition:pointed_topological_space}{map of pointed space}.
    \end{enumerate}
\end{definition}
\begin{definition} \label{definition:morphism_of_covering_spaces_of_a_topological_space}
    \begin{enumerate}
        \item Let $X$ be a \CrefAndHyperrefIfExist{definition:category_of_topological_spaces}{topological space}. 

        A \hldef{morphism of covering spaces of $X$ from $p_1: \tilde{X}_1 \to X$ to $p_2: \tilde{X}_2 \to X$} is a \CrefAndHyperrefIfExist{definition:continuous_map_of_topological_spaces}{continuous map} $f: \tilde{X}_1 \to \tilde{X}_2$ such that $p_2 \circ f = p_1$.

        \item Let $(X, x_0)$ be a \CrefAndHyperrefIfExist{definition:pointed_topological_space}{pointed topological space}. Let $p_1: (\tilde{X}_1, \tilde{x}_1) \to (X, x_0)$ to $p_2: (\tilde{X}_2, \tilde{x}_2) \to (X, x_0)$ be \CrefAndHyperrefIfExist{definition:covering_space_of_a_topological_space}{pointed coverings} of $(X, x_0)$. 
        A \hldef{morphism $f: p_1 \to p_2$ of the pointed coverings} is a pointed map $f: (\tilde{X}_1, \tilde{x}_1) \to (\tilde{X}_2,x_2)$ such that $p_2 \circ f = p_1$ and $f(\tilde{x}_1) = \tilde{x}_2$.
    \end{enumerate}
\end{definition}
\begin{definition} \label{definition:deck_transformation_of_a_covering_space}
Let $p: \tilde{X} \to X$ be a \CrefAndHyperrefIfExist{definition:covering_space_of_a_topological_space}{covering space} of a \CrefAndHyperrefIfExist{definition:category_of_topological_spaces}{topological space} $X$. 
A \hldef{deck transformation} (or \hldef{covering transformation}, or \hldef{automorphism}) of $p$ is a \CrefAndHyperrefIfExist{definition:homeomorphism_of_topological_spaces}{homeomorphism} $f: \tilde{X} \to \tilde{X}$ such that $p \circ f = p$. The set of all deck transformations forms a group under composition, denoted \hl{$\text{Aut}(p)$} or \hl{$\text{Deck}(p)$}.
\end{definition}
\begin{definition} \label{definition:universal_cover_of_a_topological_space}
Let $X$ be a \CrefAndHyperrefIfExist{definition:category_of_topological_spaces}{topological space}. 
A \hldef{universal cover of $X$} is a pair \hl{$(\tilde{X}, p)$} where $p: \tilde{X} \to X$ is a \CrefAndHyperrefIfExist{definition:covering_space_of_a_topological_space}{covering space} such that $\tilde{X}$ is \CrefAndHyperrefIfExist{definition:simply_connected_space}{simply connected}.
\end{definition}
\begin{theorem} \label{theorem:map_of_fundamental_groups_induced_by_covering_of_a_path_connected_and_locally_path_connected_space_is_injective}
Let $p: (\tilde{X}, \tilde{x}_0) \to (X, x_0)$ be a \CrefAndHyperrefIfExist{definition:covering_space_of_a_topological_space}{covering space} of a \CrefAndHyperrefIfExist{definition:path_connected_topological_space}{path-connected} and \CrefAndHyperrefIfExist{definition:locally_path_connected_topological_space}{locally path-connected space} $X$. 
The induced map on \CrefAndHyperrefIfExist{definition:homotopy_groups_of_a_pointed_topological_space}{fundamental groups} $p_*: \pi_1(\tilde{X}, \tilde{x}_0) \to \pi_1(X, x_0)$ is an \CrefAndHyperrefIfExist{definition:injective_and_projective_objects_in_a_category}{injective} \CrefAndHyperrefIfExist{definition:group_homomorphism}{group homomorphism}.
\end{theorem}
\begin{definition} \label{definition:lift_of_a_continuous_map_under_a_covering_of_topological_spaces}
Let $p: \tilde{X} \to X$ be a \CrefAndHyperrefIfExist{definition:covering_space_of_a_topological_space}{covering space} of a \CrefAndHyperrefIfExist{definition:category_of_topological_spaces}{topological space} $X$. Let $f: Y \to X$ be a \CrefAndHyperrefIfExist{definition:continuous_map_of_topological_spaces}{continuous map} of topological spaces.

A \hldef{lift of $f$ (via $p$)} is a continuous map $\tilde{f}: Y \to \tilde{X}$ satisfying $p \circ \tilde{f} = f$. 

Similarly, for a pointed covering $p:(\tilde{X}, x_0) \to (X,x_0)$ and a \CrefAndHyperrefIfExist{definition:pointed_topological_space}{pointed morphism} $f:(Y,y_0) \to (x,x_0)$, a \hldef{lift of $f$ (via $p$)} is a pointed map $\tilde{f}: (Y,y_0) \to (\tilde{X},x_0)$ satisfying $p \circ \tilde{f} = f$.
\end{definition}

\begin{definition} \label{definition:normal_covering_space_of_a_topological_space}
A \CrefAndHyperrefIfExist{definition:covering_space_of_a_topological_space}{covering space} $p: \tilde{X} \to X$ of a \CrefAndHyperrefIfExist{definition:category_of_topological_spaces}{topological space} $X$ is \hldef{normal} (or \hldef{Galois)} if for every $\tilde{x}_0 \in X_0$, the image of the induced map $p_*: \pi_1(\tilde{X},\tilde{x}_0) \to \pi_1(\tilde{X},p(\tilde{x}_0))$ on \CrefAndHyperrefIfExist{definition:homotopy_groups_of_a_pointed_topological_space}{fundamental groups} is the normal subgroup of $\pi_1(\tilde{X},p(\tilde{x}_0))$ 

% the \CrefAndHyperrefIfExist{definition:group}{group} of deck transformations $\text{Aut}(p)$ acts transitively on the \CrefAndHyperrefIfExist{definition:fiber_of_a_map_of_topological_spaces_over_a_point}{fiber} $p^{-1}(x)$.
% \TODO{transitive action}
\end{definition}
\begin{theorem} \label{theorem:lift_exists_for_a_map_of_pointed_spaces_under_a_covering_of_pointed_spaces_iff_images_of_fundamental_groups_are_contained}
Let $p: (\tilde{X}, \tilde{x}_0) \to (X, x_0)$ be a \CrefAndHyperrefIfExist{definition:covering_space_of_a_topological_space}{covering space} and let $f: (Y, y_0) \to (X, x_0)$ be a \CrefAndHyperrefIfExist{definition:continuous_map_of_topological_spaces}{continuous map} where $Y$ is \CrefAndHyperrefIfExist{definition:path_connected_topological_space}{path-connected} and \CrefAndHyperrefIfExist{definition:locally_path_connected_topological_space}{locally path-connected}. 
A lift $\tilde{f}: (Y, y_0) \to (\tilde{X}, \tilde{x}_0)$ exists if and only if the image of the \CrefAndHyperrefIfExist{definition:homotopy_groups_of_a_pointed_topological_space}{fundamental group} of $(Y,y_0)$ is contained in the image of the fundamental group of $(\tilde{X}, \tilde{x}_0)$, i.e. 
$$f_*(\pi_1(Y, y_0)) \subseteq p_*(\pi_1(\tilde{X}, \tilde{x}_0))$$
If such a lift exists, it is unique.

\TextIfExists{theorem:lift_for_a_map_of_pointed_spaces_to_a_map_of_covers_exists_iff_images_of_fundamental_groups_are_contained}{
See also \Cref{theorem:lift_for_a_map_of_pointed_spaces_to_a_map_of_covers_exists_iff_images_of_fundamental_groups_are_contained}.
}
\end{theorem}
\begin{theorem} \label{theorem:lift_for_a_map_of_pointed_spaces_to_a_map_of_covers_exists_iff_images_of_fundamental_groups_are_contained}
Let $p_X: (\tilde{X}, \tilde{x}_0) \to (X, x_0)$ and $p_Y: (\tilde{Y}, \tilde{y}_0) \to (Y, y_0)$ be covering spaces of \CrefAndHyperrefIfExist{definition:pointed_topological_space}{pointed topological spaces}. Let $f: (X, x_0) \to (Y, y_0)$ be a \CrefAndHyperrefIfExist{definition:continuous_map_of_topological_spaces}{continuous map}, and assume $\tilde{X}$ is \CrefAndHyperrefIfExist{definition:path_connected_topological_space}{path-connected} and \CrefAndHyperrefIfExist{definition:locally_path_connected_topological_space}{locally path-connected}. 

There exists a unique continuous lift $\tilde{f}: (\tilde{X}, \tilde{x}_0) \to (\tilde{Y}, \tilde{y}_0)$ such that $p_Y \circ \tilde{f} = f \circ p_X$ if and only if 
$$f_*(p_{X*}(\pi_1(\tilde{X}, \tilde{x}_0))) \subseteq p_{Y*}(\pi_1(\tilde{Y}, \tilde{y}_0))$$
In the specific case where $f$ is the identity on $X$, this provides the necessary and sufficient condition for the existence of a morphism of covering spaces (pointed or unpointed). Any such morphism between connected covering spaces is itself a covering map.
\TextIfExists{theorem:lift_exists_for_a_map_of_pointed_spaces_under_a_covering_of_pointed_spaces_iff_images_of_fundamental_groups_are_contained}{
See also \Cref{theorem:lift_exists_for_a_map_of_pointed_spaces_under_a_covering_of_pointed_spaces_iff_images_of_fundamental_groups_are_contained}.
}
\end{theorem}
\begin{theorem} \label{theorem:universal_cover_exists_for_connected_locally_path_connected_semi_locally_simply_connected_topological_space}
Let $X$ be a \CrefAndHyperrefIfExist{definition:disconnected_and_connected_topological_spaces}{connected}, \CrefAndHyperrefIfExist{definition:locally_path_connected_topological_space}{locally path-connected}, and \CrefAndHyperrefIfExist{definition:semi_locally_simply_connected_space}{semi-locally simply connected} \CrefAndHyperrefIfExist{definition:A1_homotopy_theoretic_space_over_a_noetherian_scheme_of_finite_dimension}{space} with \CrefAndHyperrefIfExist{definition:basis_for_a_topology}{base} point $x_0$. The \CrefAndHyperrefIfExist{definition:universal_cover_of_a_topological_space}{universal cover} $\tilde{X}$ can be constructed as the space of \CrefAndHyperrefIfExist{definition:homotopy_class_of_maps_of_topological_spaces_relative_to_a_subset}{homotopy classes} of \CrefAndHyperrefIfExist{definition:path_and_loop_in_a_topological_space}{paths} starting at $x_0$:
$$\tilde{X} = \{ [\gamma] \mid \gamma: [0,1] \to X, \gamma(0) = x_0 \}$$
with the covering map $p: \tilde{X} \to X$ given by $p([\gamma]) = \gamma(1)$. For any connected \CrefAndHyperrefIfExist{definition:covering_space_of_a_topological_space}{covering space} $q: Y \to X$, there exists a covering map $r: \tilde{X} \to Y$ such that $q \circ r = p$.
\end{theorem}
\begin{proposition} \label{proposition:isom_class_of_conn_pointed_covers_of_a_conn_loc_path_conn_semiloc_simp_conn_pointed_space_corresponds_to_subgroup_of_fund_gp}
Let $X$ be a \CrefAndHyperrefIfExist{definition:disconnected_and_connected_topological_spaces}{connected}, \CrefAndHyperrefIfExist{definition:locally_path_connected_topological_space}{locally path-connected}, and \CrefAndHyperrefIfExist{definition:semi_locally_simply_connected_space}{semi-locally simply connected} space with \CrefAndHyperrefIfExist{definition:pointed_topological_space}{base point} $x_0$. There are \CrefAndHyperrefIfExist{definition:injective_surjective_bijective_map_of_sets}{bijective} correspondences:
\begin{enumerate}
    \item Between the set of isomorphism classes of \CrefAndHyperrefIfExist{definition:disconnected_and_connected_topological_spaces}{connected} \CrefAndHyperrefIfExist{definition:covering_space_of_a_topological_space}{pointed covering spaces} $p: (\tilde{X}, \tilde{x}_0) \to (X, x_0)$ and the set of \CrefAndHyperrefIfExist{definition:subgroup_of_a_group}{subgroups} of $\pi_1(X, x_0)$, via $((\tilde{X}, \tilde{x}_0), p) \mapsto p_*(\pi_1(\tilde{X}, \tilde{x}_0))$.
    \item Between the set of isomorphism classes of connected covering spaces $p: \tilde{X} \to X$ and the set of conjugacy classes of subgroups of $\pi_1(X, x_0)$, via $p \mapsto [p_*(\pi_1(\tilde{X}, \tilde{x}_0))]$.
    \TODO{conjugacy class }
\end{enumerate}
\end{proposition}
\begin{theorem} \label{theorem:deck_transformation_group_of_a_connected_pointed_covering_is_quotient}
Let $p: (\tilde{X}, \tilde{x}_0) \to (X, x_0)$ be a covering with $\tilde{X}$ (and hence $X$) \CrefAndHyperrefIfExist{definition:disconnected_and_connected_topological_spaces}{connected}. 
The \CrefAndHyperrefIfExist{definition:deck_transformation_of_a_covering_space}{group of deck transformations} $\text{Aut}(p)$ is isomorphic to the \CrefAndHyperrefIfExist{definition:quotient_of_a_group_by_a_normal_subgroup}{quotient} $N(H)/H$, where $H = p_*(\pi_1(\tilde{X}, \tilde{x}_0))$ and $N(H)$ is its \CrefAndHyperrefIfExist{definition:normalizer_of_a_subgroup_of_a_group}{normalizer} in $\pi_1(X, x_0)$. If the covering is \CrefAndHyperrefIfExist{definition:normal_covering_space_of_a_topological_space}{normal} (i.e., $H$ is a \CrefAndHyperrefIfExist{definition:normal_subgroup_and_characteristic_subgroup_of_an_algebraic_group_over_a_scheme}{normal subgroup} of $\pi_1(X, x_0)$), then:
$$\text{Aut}(p) \cong \pi_1(X, x_0) / p_*(\pi_1(\tilde{X}, \tilde{x}_0))$$
For the universal cover, $\text{Aut}(p) \cong \pi_1(X, x_0)$.
\end{theorem}


\section{Miscellaneous definitions}
\begin{definition} \label{definition:conjugacy_invariant_subset_of_a_group}
    Let $G$ be a group. A subset $S \subset G$ is said to be \hldef{conjugacy invariant} or \hldef{invariant under conjugation by $G$} if for all $g \in G$ and all $s \in G$, the \CrefAndHyperrefIfExist{definition:conjugation_of_group_elements_and_subgroup_by_a_group_element}{conjugate} $gSg^{-1}$ is an element of $S$. 
\end{definition}
